\documentclass[11pt,a4paper]{article}

% ─── Packages ────────────────────────────────────────────────────────────
\usepackage[utf8]{inputenc}
\usepackage[T1]{fontenc}
\usepackage{amsmath,amssymb,amsthm,mathtools}
\usepackage{geometry}
\usepackage{hyperref}
\usepackage{cleveref}
\usepackage{algorithm}
\usepackage{algpseudocode}
\usepackage{booktabs}
\usepackage{graphicx}
\usepackage{tikz}
\usepackage{tikz-cd}
\usetikzlibrary{positioning,arrows.meta,decorations.pathreplacing,calc,shapes.geometric}
\usepackage{enumitem}
\usepackage{caption}
\usepackage{subcaption}
\usepackage{thmtools}
\usepackage{xcolor}
\usepackage{fancyhdr}
\usepackage{titlesec}

\geometry{margin=1in}

% ─── Theorem environments ────────────────────────────────────────────────
\theoremstyle{definition}
\newtheorem{definition}{Definition}[section]
\newtheorem{example}[definition]{Example}

\theoremstyle{plain}
\newtheorem{theorem}[definition]{Theorem}
\newtheorem{proposition}[definition]{Proposition}
\newtheorem{lemma}[definition]{Lemma}
\newtheorem{corollary}[definition]{Corollary}

\theoremstyle{remark}
\newtheorem{remark}[definition]{Remark}

% ─── Custom commands ─────────────────────────────────────────────────────
\newcommand{\R}{\mathbb{R}}
\newcommand{\N}{\mathbb{N}}
\newcommand{\E}{\mathbb{E}}
\newcommand{\Hcal}{\mathcal{H}}
\newcommand{\Scal}{\mathcal{S}}
\newcommand{\Lcal}{\mathcal{L}}
\newcommand{\Mcal}{\mathcal{M}}
\newcommand{\Tcal}{\mathcal{T}}
\newcommand{\Ocal}{\mathcal{O}}
\newcommand{\Vcal}{\mathcal{V}}
\newcommand{\softmax}{\operatorname{softmax}}
\newcommand{\sigmoid}{\sigma}
\newcommand{\diag}{\operatorname{diag}}
\newcommand{\tr}{\operatorname{tr}}
\newcommand{\rank}{\operatorname{rank}}
\newcommand{\Rot}{\operatorname{Rot}}
\newcommand{\Ent}{\operatorname{H}}
\DeclareMathOperator*{\argmin}{arg\,min}
\DeclareMathOperator*{\argmax}{arg\,max}

% ─── Header/footer ──────────────────────────────────────────────────────
\pagestyle{fancy}
\fancyhf{}
\fancyhead[L]{\textsc{Felix-LM}}
\fancyhead[R]{\textsc{Design Document v0.1}}
\fancyfoot[C]{\thepage}
\renewcommand{\headrulewidth}{0.4pt}

% ─── Title ───────────────────────────────────────────────────────────────
\title{%
  \textbf{Felix-LM: Multi-Stream Progressive Merging\\
  for Causal Language Modeling}\\[0.5em]
  \large Architecture Design Document \& Theoretical Foundation
}
\author{%
  Caleb Gross\\
  {\small\textit{[Affiliation]}}\\
  {\small\texttt{caleb@appsprout.dev}}
}
\date{Draft --- \today}

\begin{document}
\maketitle
\thispagestyle{fancy}

% ═════════════════════════════════════════════════════════════════════════
\begin{abstract}
We propose \textbf{Felix-LM}, a novel causal language model architecture 
derived from the computational principles of the Felix multi-agent 
framework. Where standard transformer architectures process representations 
through a uniform stack of identical layers, Felix-LM introduces 
\emph{Multi-Stream Progressive Merging} (MSPM): the input is processed by 
multiple independent computational streams that progressively merge into a 
single output stream through a sequence of structured convergence operations.

We formalize the claim that this architecture induces a \emph{helical funnel 
trajectory} in representation space---a path characterized by simultaneous 
rotational mixing and radial compression as representations move through the 
network's depth. We provide the mathematical framework for this trajectory, 
specify the architecture precisely, analyze its gradient flow and 
information-theoretic properties, and propose an ablation study to test 
whether the multi-stream progressive merging inductive bias improves 
perplexity on standard language modeling benchmarks relative to 
parameter-matched uniform transformers.
\end{abstract}

\tableofcontents
\newpage

% ═════════════════════════════════════════════════════════════════════════
\section{Introduction}\label{sec:intro}

\subsection{Motivation}

The dominant architecture for autoregressive language modeling is the 
decoder-only transformer \cite{vaswani2017attention, radford2019language}, 
consisting of $L$ structurally identical layers, each comprising multi-head 
self-attention and a position-wise feedforward network. Every layer operates 
on a representation of fixed dimension $d_{\text{model}}$, and the output of 
each layer is added to the input via a residual connection. We refer to this 
as the \emph{uniform} architecture.

Empirical evidence suggests that this uniformity is suboptimal: early layers 
in trained transformers predominantly capture local syntactic features, while 
later layers encode global semantic and reasoning structure 
\cite{tenney2019bert, clark2019what}. The uniform architecture allocates 
identical computational resources to both regimes, suggesting that a 
\emph{heterogeneous} architecture---one whose structure varies with 
depth---might achieve better performance per parameter.

\subsection{The Felix Framework}

The Felix multi-agent framework\footnote{Developed by C.\ [Surname]; 
codebase and documentation available at [repository].} provides the 
conceptual foundation for this work. Felix models a team of AI agents 
traversing a helical path in 3D space, governed by the parametric equations:
\begin{align}
  x(t) &= R(t)\cos\bigl(2\pi n t\bigr), \label{eq:felix-x}\\
  y(t) &= R(t)\sin\bigl(2\pi n t\bigr), \label{eq:felix-y}\\
  z(t) &= H(1 - t), \label{eq:felix-z}
\end{align}
where $t \in [0,1]$ is the progression parameter, $n$ is the number of 
helical turns, $H$ is the total height, and the radius follows an 
exponential tapering function:
\begin{equation}\label{eq:felix-radius}
  R(t) = R_{\min} \left(\frac{R_{\max}}{R_{\min}}\right)^{1-t}.
\end{equation}

In Felix, agents at the top of the helix ($t \approx 0$) perform broad, 
divergent exploration; agents at the bottom ($t \approx 1$) perform focused, 
convergent synthesis. A central hub (the \emph{CentralPost}) sits on the 
helix's axis and performs final integration of all agent outputs. The key 
computational principles are:
\begin{enumerate}[label=(\roman*)]
  \item \textbf{Parallel exploration, progressive convergence}: Multiple 
    independent processing paths narrow to a single output.
  \item \textbf{Depth-dependent computation type}: Early processing is 
    broad and divergent; late processing is focused and convergent.
  \item \textbf{Confidence-gated resource allocation}: Computation is 
    allocated adaptively based on a confidence signal.
  \item \textbf{Structured communication at convergence points}: 
    Cross-stream information exchange occurs at specific merge events, not 
    globally at every step.
\end{enumerate}

\subsection{Thesis Statement}

We propose that these four principles can be instantiated as an inductive 
bias in the architecture of a causal language model, yielding a model 
that traces a \emph{helical funnel} in representation space: a trajectory 
characterized by monotonic reduction in the number of independent 
processing streams (the ``funnel'') combined with learned rotational mixing 
at stream-merge boundaries (the ``helix''). We hypothesize that this 
inductive bias improves language modeling perplexity relative to a 
parameter-matched uniform transformer, particularly on sequences requiring 
multi-step reasoning.

% ═════════════════════════════════════════════════════════════════════════
\section{Mathematical Preliminaries}\label{sec:prelim}

\subsection{Notation}

Let $\Vcal$ denote a vocabulary of size $V$. An input sequence is 
$\mathbf{x} = (x_1, \ldots, x_T)$ with $x_i \in \Vcal$. A causal 
language model defines $p(x_{t+1} \mid x_1, \ldots, x_t)$ for all $t$.

We write $\R^{d}$ for $d$-dimensional Euclidean space. For a matrix 
$W \in \R^{m \times n}$, we write $W^\top$ for its transpose and 
$\|W\|_F$ for its Frobenius norm. The identity matrix is $I_d \in 
\R^{d \times d}$. We write $\sigmoid(\cdot)$ for the logistic sigmoid 
function and $\softmax(\cdot)$ for the softmax function.

\subsection{Representation Space and Trajectories}\label{sec:rep-traj}

\begin{definition}[Representation trajectory]\label{def:rep-traj}
Let $f: \R^{d_{\text{in}}} \to \R^{d_{\text{out}}}$ be a neural network 
of depth $L$ with intermediate representations $\mathbf{h}^{(0)}, 
\mathbf{h}^{(1)}, \ldots, \mathbf{h}^{(L)}$. The \emph{representation 
trajectory} of an input $\mathbf{x}$ is the sequence
\[
  \gamma(\mathbf{x}) = \bigl(\mathbf{h}^{(0)}, \mathbf{h}^{(1)}, \ldots, 
  \mathbf{h}^{(L)}\bigr),
\]
where $\mathbf{h}^{(\ell)} \in \R^{d_\ell}$ is the representation at 
depth $\ell$. When $d_\ell$ varies with $\ell$, we embed all 
representations into a common ambient space $\R^{d_{\max}}$ via 
zero-padding for the purpose of geometric analysis.
\end{definition}

\begin{definition}[Effective dimensionality]\label{def:eff-dim}
For a representation $\mathbf{h} \in \R^d$, the \emph{effective 
dimensionality} with respect to a gate vector $\mathbf{g} \in [0,1]^d$ is
\[
  d_{\text{eff}}(\mathbf{h}, \mathbf{g}) = \|\mathbf{g}\|_1 = 
  \sum_{i=1}^{d} g_i.
\]
This generalizes hard dimensionality reduction: when $g_i \in \{0,1\}$, 
$d_{\text{eff}}$ counts the number of active dimensions. For soft gates 
$g_i \in (0,1)$, it provides a smooth interpolation.
\end{definition}

\begin{definition}[Helical trajectory in representation space]
\label{def:helical-traj}
A representation trajectory $\gamma$ is \emph{helical} if there exists a 
sequence of orthogonal transformations $\{Q^{(\ell)} \in O(d_{\max})\}$ 
and a monotonically decreasing effective dimensionality sequence 
$\{d_{\text{eff}}^{(\ell)}\}$ such that the transition between successive 
representations can be decomposed as:
\begin{equation}\label{eq:helical-decomp}
  \mathbf{h}^{(\ell+1)} = Q^{(\ell)} \cdot 
  \bigl(\mathbf{g}^{(\ell)} \odot f^{(\ell)}(\mathbf{h}^{(\ell)})\bigr),
\end{equation}
where $f^{(\ell)}$ is the layer computation, $\mathbf{g}^{(\ell)}$ is a 
gating vector with $\|\mathbf{g}^{(\ell)}\|_1 \leq 
\|\mathbf{g}^{(\ell-1)}\|_1$, $\odot$ denotes elementwise multiplication, 
and $Q^{(\ell)}$ is an orthogonal rotation.

The \emph{angular velocity} of the trajectory is 
$\omega^{(\ell)} = \arccos\bigl(\frac{\tr(Q^{(\ell)})}{d_{\max}}\bigr)$, 
and the \emph{radial compression rate} is 
$\rho^{(\ell)} = d_{\text{eff}}^{(\ell+1)} / d_{\text{eff}}^{(\ell)} 
\leq 1$.
\end{definition}

\begin{remark}
In a standard uniform transformer with residual connections, the trajectory 
is approximately \emph{linear} (representations are incrementally shifted 
in a consistent direction) with constant effective dimensionality. The 
helical trajectory introduces both rotation and compression, forcing the 
network to progressively refine and consolidate its representation. This is 
the mathematical content of the Felix helical metaphor.
\end{remark}

\subsection{Rotary Position Embeddings (RoPE)}\label{sec:rope}

Rotary Position Embeddings \cite{su2021roformer} encode token position $m$ 
by applying a block-diagonal rotation matrix $\Rot(m) \in \R^{d \times d}$ 
to query and key vectors before computing attention. For dimension pair 
$(2i, 2i+1)$ with frequency $\theta_i = 10000^{-2i/d}$:
\begin{equation}\label{eq:rope}
  \Rot(m)_{2i:2i+1} = 
  \begin{pmatrix} 
    \cos(m\theta_i) & -\sin(m\theta_i) \\ 
    \sin(m\theta_i) & \phantom{-}\cos(m\theta_i) 
  \end{pmatrix}.
\end{equation}
The key property is that the inner product between rotated queries and keys 
depends only on relative position: 
$\langle \Rot(m)\mathbf{q}, \Rot(n)\mathbf{k}\rangle = 
\langle \Rot(m-n)\mathbf{q}, \mathbf{k}\rangle$.

% ═════════════════════════════════════════════════════════════════════════
\section{Architecture}\label{sec:arch}

Felix-LM is organized into $K$ \emph{stages}, indexed $k = 0, 1, \ldots, 
K-1$. Stage $k$ contains $S_k$ independent processing \emph{streams}, each 
of model dimension $d_k$. The number of streams decreases monotonically: 
$S_0 > S_1 > \cdots > S_{K-1} = 1$. Between consecutive stages, a 
\emph{merge operation} reduces the stream count by combining pairs of 
streams. Each stream within a stage consists of $L_k$ transformer layers 
(self-attention + FFN), with architecture details that may vary by stage.

The overall computation proceeds:
\begin{equation}\label{eq:overall}
  \underbrace{\text{Embed}}_{\text{input}} 
  \;\to\; 
  \underbrace{S_0 \text{ streams} \times L_0 \text{ layers}}_{\text{Stage 0}}
  \;\xrightarrow{\text{Merge}}\;
  \underbrace{S_1 \text{ streams} \times L_1 \text{ layers}}_{\text{Stage 1}}
  \;\xrightarrow{\text{Merge}}\;
  \cdots
  \;\xrightarrow{\text{Merge}}\;
  \underbrace{1 \text{ stream} \times L_{K-1} \text{ layers}}_{\text{Stage } K{-}1}
  \;\to\;
  \underbrace{\text{Head}}_{\text{output}}.
\end{equation}

\subsection{Input Embedding and Stream Initialization}\label{sec:embed}

Given an input sequence $\mathbf{x} = (x_1, \ldots, x_T)$, we compute a 
shared token embedding:
\begin{equation}
  \mathbf{e}_t = \text{Embed}(x_t) \in \R^{d_{\text{embed}}},
  \quad t = 1, \ldots, T.
\end{equation}

Each stream $s \in \{1, \ldots, S_0\}$ in Stage~0 receives the input 
through a learned \emph{stream projection}:
\begin{equation}\label{eq:stream-init}
  \mathbf{h}_{t,s}^{(0,0)} = W_s^{\text{init}} \mathbf{e}_t + 
  \mathbf{b}_s^{\text{init}}, 
  \quad W_s^{\text{init}} \in \R^{d_0 \times d_{\text{embed}}},
\end{equation}
where each projection $W_s^{\text{init}}$ is initialized to encourage 
diverse initial representations across streams (see \S\ref{sec:init}).

\begin{remark}
Each stream can be understood as attending to a different ``aspect'' of 
the input, analogous to how Felix spawns agents of different types 
(Research, Analysis, Critic) that examine the same task from different 
perspectives. The stream projections play the role of role specialization.
\end{remark}

\subsection{Intra-Stream Computation}\label{sec:intra-stream}

Within stage $k$, each stream $s$ independently processes its 
representation through $L_k$ transformer layers. We denote the 
representation at layer $\ell$ within stage $k$, stream $s$, token 
position $t$ as $\mathbf{h}_{t,s}^{(k,\ell)} \in \R^{d_k}$.

Each layer consists of:
\begin{align}
  \tilde{\mathbf{h}} &= \mathbf{h}_{t,s}^{(k,\ell)} + 
    \text{Attn}_{k,\ell,s}\bigl(
      \text{LN}(\mathbf{h}_{:,s}^{(k,\ell)})\bigr)_t, 
    \label{eq:attn-residual} \\
  \mathbf{h}_{t,s}^{(k,\ell+1)} &= \tilde{\mathbf{h}} + 
    \text{FFN}_{k,\ell,s}\bigl(\text{LN}(\tilde{\mathbf{h}})\bigr),
    \label{eq:ffn-residual}
\end{align}
where $\text{LN}$ denotes RMSNorm \cite{zhang2019root}, and the 
attention and FFN modules are specific to stage $k$, layer $\ell$, and 
stream $s$. The notation $\mathbf{h}_{:,s}^{(k,\ell)}$ denotes the full 
sequence of representations for stream $s$ (used as keys/values in 
attention).

\paragraph{Heterogeneous attention by stage.}
A key design principle is that the \emph{type} of attention varies with 
stage, reflecting Felix's principle of depth-dependent computation:

\begin{center}
\begin{tabular}{@{}clll@{}}
\toprule
\textbf{Stage} & \textbf{Felix Analog} & \textbf{Attention Type} & 
  \textbf{Rationale} \\
\midrule
$k = 0$ & Research (exploration) & Linear / SSM & 
  $\Ocal(T)$: broad, cheap context \\
$k = 1, \ldots, K{-}2$ & Analysis (convergence) & Sliding window & 
  $\Ocal(Tw)$: local precision \\
$k = K{-}1$ & CentralPost (synthesis) & Full causal & 
  $\Ocal(T^2)$: maximum precision \\
\bottomrule
\end{tabular}
\end{center}

This allocation concentrates quadratic-cost attention at the final stage, 
where only a single stream operates, and uses efficient linear-time 
mechanisms in the early stages where multiple parallel streams would 
otherwise multiply the cost.

\paragraph{Weight sharing.}
Streams within the same stage may share weights (all streams identical, 
differentiated only by their initial projection) or have independent 
weights (each stream is a distinct sub-network). We denote these 
configurations as \textsc{Shared} and \textsc{Independent}, respectively, 
and treat this as an ablation variable.

\subsection{Merge Operation}\label{sec:merge}

The merge operation is the architectural core of Felix-LM, implementing 
the progressive convergence that defines the helical funnel.

\begin{definition}[Stream merge]\label{def:merge}
Let streams $s$ and $s'$ at the end of stage $k$ have representations 
$\mathbf{h}_{t,s}^{(k, L_k)}, \mathbf{h}_{t,s'}^{(k, L_k)} \in 
\R^{d_k}$ for each token $t$. The \emph{merge operation} produces a single 
stream representation $\mathbf{m}_t \in \R^{d_{k+1}}$ for the next stage:
\begin{equation}\label{eq:merge}
  \mathbf{m}_t = \underbrace{P_{k}}_{\text{projection}} 
  \Bigl( 
    \underbrace{\mathbf{g}_{k}([\mathbf{h}_{t,s}; \mathbf{h}_{t,s'}])}_%
      {\text{gated combination}} 
  \Bigr),
\end{equation}
where $[\cdot\,;\,\cdot]$ denotes concatenation along the feature axis, 
$P_k \in \R^{d_{k+1} \times 2d_k}$ is a learned projection, and 
$\mathbf{g}_k: \R^{2d_k} \to [0,1]^{2d_k}$ is a gating function:
\begin{equation}\label{eq:gate}
  \mathbf{g}_k(\mathbf{z}) = \sigmoid\bigl(W_k^g \mathbf{z} + 
  \mathbf{b}_k^g\bigr), 
  \quad W_k^g \in \R^{2d_k \times 2d_k}.
\end{equation}
The full merge operation is thus:
\begin{equation}\label{eq:merge-full}
  \mathbf{m}_t = P_k \Bigl( \sigmoid\bigl(W_k^g 
  [\mathbf{h}_{t,s}; \mathbf{h}_{t,s'}] + \mathbf{b}_k^g\bigr) \odot 
  [\mathbf{h}_{t,s}; \mathbf{h}_{t,s'}] \Bigr).
\end{equation}
\end{definition}

\begin{remark}[Connection to Felix's CentralPost]
In Felix, the CentralPost synthesizes agent outputs with adaptive 
parameters (temperature $\in [0.2, 0.4]$, token budget $\in 
[1500, 3000]$). The merge gate $\mathbf{g}_k$ plays an analogous role: it 
learns which features from each stream to preserve (high gate value) versus 
suppress (low gate value) during synthesis. The projection $P_k$ 
corresponds to the CentralPost's final synthesis prompt.
\end{remark}

\subsubsection{Cross-Stream Attention at Merge Points}
\label{sec:cross-attn}

Before applying the gated projection, we optionally perform a single layer 
of \emph{cross-stream attention}, where each stream attends to the other:
\begin{align}
  \hat{\mathbf{h}}_{t,s} &= \mathbf{h}_{t,s}^{(k,L_k)} + 
    \text{CrossAttn}\bigl(\mathbf{h}_{t,s}^{(k,L_k)},\; 
    \mathbf{h}_{:,s'}^{(k,L_k)}\bigr), \label{eq:cross-attn}
\end{align}
where $\text{CrossAttn}(Q\text{-source}, KV\text{-source})$ is standard 
multi-head attention with queries from one stream and keys/values from the 
other. This allows streams to exchange information before fusion, 
implementing Felix's principle of structured communication at convergence 
points. We treat the inclusion of cross-stream attention as an ablation 
variable.

\subsubsection{The Merge Projection as Rotation}
\label{sec:merge-rotation}

We can decompose the merge projection $P_k$ via polar decomposition:
\begin{equation}\label{eq:polar}
  P_k = U_k \Sigma_k V_k^\top, 
\end{equation}
where $U_k \in \R^{d_{k+1} \times d_{k+1}}$ and $V_k \in \R^{2d_k 
\times 2d_k}$ are orthogonal, and $\Sigma_k$ is the singular value matrix. 
The matrix $V_k^\top$ rotates the concatenated representation in 
$\R^{2d_k}$, $\Sigma_k$ scales (compresses) dimensions, and $U_k$ rotates 
the result in $\R^{d_{k+1}}$.

This decomposition reveals that every merge naturally performs 
\emph{rotation followed by compression followed by rotation}---the 
mathematical structure of a helical step. We may optionally constrain 
$P_k$ to have orthogonal rows ($P_k P_k^\top = I_{d_{k+1}}$) to 
guarantee \emph{isometric} compression: the merge preserves inner products 
and thus similarity structure, ensuring that information is consolidated 
without distortion. This constraint can be enforced via Cayley 
parameterization:
\begin{equation}\label{eq:cayley}
  P_k = (I - A_k)(I + A_k)^{-1} \cdot \Pi,
\end{equation}
where $A_k$ is a learnable skew-symmetric matrix and $\Pi \in 
\R^{d_{k+1} \times 2d_k}$ is a fixed truncation matrix. We treat this 
orthogonality constraint as an ablation variable.

\subsection{Depth-Extended Rotary Embeddings}\label{sec:depth-rope}

Standard RoPE (\S\ref{sec:rope}) encodes token position via rotation. 
We extend this to also encode \emph{depth} (layer index), providing the 
``twist'' of the helical trajectory.

\begin{definition}[Depth-extended RoPE]\label{def:depth-rope}
At global layer index $\ell_{\text{global}}$ (counting across all stages) 
and token position $m$, the rotation applied to dimension pair 
$(2i, 2i+1)$ of queries and keys is:
\begin{equation}\label{eq:depth-rope}
  \Rot(m, \ell_{\text{global}})_{2i:2i+1} = 
  \begin{pmatrix} 
    \cos\phi_i & -\sin\phi_i \\ 
    \sin\phi_i & \phantom{-}\cos\phi_i 
  \end{pmatrix},
\end{equation}
where the composite angle is:
\begin{equation}\label{eq:composite-angle}
  \phi_i(m, \ell_{\text{global}}) = m \cdot \theta_i^{\text{pos}} + 
  \ell_{\text{global}} \cdot \theta_i^{\text{depth}},
\end{equation}
with position frequencies $\theta_i^{\text{pos}} = 
\beta_{\text{pos}}^{-2i/d}$ and depth frequencies 
$\theta_i^{\text{depth}} = \alpha_i \cdot 
\frac{2\pi n}{L_{\text{total}}}$.
\end{definition}

Here $n$ is the number of helical turns (directly from the Felix 
configuration), $L_{\text{total}} = \sum_k L_k$ is the total layer count, 
$\beta_{\text{pos}}$ is the RoPE base frequency (typically $10000$), and 
$\alpha_i$ is a per-dimension scaling factor that controls the coupling 
strength between depth and position. Setting $\alpha_i = 0$ recovers 
standard RoPE; setting $\alpha_i = 1$ applies the full helical twist.

\begin{proposition}[Relative depth encoding]\label{prop:rel-depth}
Under depth-extended RoPE, the attention score between query at position 
$m$, depth $\ell$ and key at position $m'$, depth $\ell'$ in dimension 
pair $i$ depends only on the relative position $\Delta m = m - m'$ and 
relative depth $\Delta\ell = \ell - \ell'$:
\begin{equation}
  \langle \Rot(m, \ell)\mathbf{q}_i, \Rot(m', \ell')\mathbf{k}_i\rangle 
  = \langle \Rot(\Delta m, \Delta\ell)\mathbf{q}_i, 
  \mathbf{k}_i\rangle.
\end{equation}
\end{proposition}
\begin{proof}
This follows from the group homomorphism property of rotation matrices: 
$\Rot(\phi_1)\Rot(\phi_2) = \Rot(\phi_1 + \phi_2)$, since $\phi_i$ is 
linear in both $m$ and $\ell_{\text{global}}$.
\end{proof}

\begin{remark}
In causal (decoder-only) language modeling, cross-depth attention does not 
occur---each layer only attends within itself. The depth-extended RoPE 
nonetheless affects the geometry: it ensures that representations at 
successive depths exist in \emph{rotated} coordinate frames, which 
influences how the residual stream accumulates information and how the 
merge projections interact with different depths' representations. The 
relative-depth property becomes directly relevant if skip connections 
across stages are introduced (\S\ref{sec:extensions}).
\end{remark}

\subsection{Confidence-Gated Early Exit}\label{sec:early-exit}

Felix agents with high confidence accelerate through the helix. The LLM 
analog is that tokens whose next-token prediction is ``easy'' need not 
traverse all $K$ stages.

At each merge boundary (after stage $k$, $k < K-1$), we attach a 
lightweight \emph{exit head}:
\begin{equation}\label{eq:exit-head}
  \hat{p}_k(x_{t+1} \mid \mathbf{x}_{\leq t}) = 
  \softmax\bigl(W_k^{\text{exit}} \cdot 
  \text{LN}(\bar{\mathbf{h}}_t^{(k)}) + \mathbf{b}_k^{\text{exit}}\bigr),
\end{equation}
where $\bar{\mathbf{h}}_t^{(k)}$ is the merged representation at the end 
of stage $k$ (or the mean over streams if exit is evaluated before merge).

\subsubsection{Cross-Stream Agreement as Confidence}

When multiple streams are active, we define a \emph{cross-stream 
agreement} score:
\begin{equation}\label{eq:agreement}
  \text{Agree}_t^{(k)} = \frac{1}{\binom{S_k}{2}} 
  \sum_{s < s'} \cos\bigl(\mathbf{h}_{t,s}^{(k, L_k)},\; 
  \mathbf{h}_{t,s'}^{(k, L_k)}\bigr),
\end{equation}
where $\cos(\cdot, \cdot)$ denotes cosine similarity. This measures the 
degree to which independent processing streams have converged to the same 
representation for token $t$.

\begin{definition}[Early exit criterion]\label{def:exit}
Token $t$ exits at stage $k$ if:
\begin{equation}\label{eq:exit-crit}
  \text{Agree}_t^{(k)} > \tau_{\text{conf}} 
  \quad\text{and}\quad 
  \Ent\bigl[\hat{p}_k(\cdot \mid \mathbf{x}_{\leq t})\bigr] < 
  \tau_{\text{ent}},
\end{equation}
where $\Ent[\cdot]$ denotes entropy and $\tau_{\text{conf}}, 
\tau_{\text{ent}}$ are thresholds. The first condition requires 
cross-stream consensus (Felix's convergence detection); the second 
requires the prediction itself to be low-entropy.
\end{definition}

\subsubsection{Training Objective with Deep Supervision}

During training, all tokens traverse all stages (no early exit). The loss 
combines predictions from all exit points and the final output:
\begin{equation}\label{eq:loss}
  \Lcal = \sum_{k=0}^{K-1} \lambda_k \cdot \Lcal_k, 
  \quad\text{where}\quad 
  \Lcal_k = -\frac{1}{T}\sum_{t=1}^{T} 
  \log \hat{p}_k(x_{t+1} \mid \mathbf{x}_{\leq t}),
\end{equation}
with weights $\lambda_k$ satisfying $\sum_k \lambda_k = 1$ and 
$\lambda_{K-1} > \lambda_k$ for $k < K-1$ (the final stage is weighted 
most heavily). This deep supervision encourages useful representations at 
every stage, enabling meaningful early exit at inference time.

\subsection{Output Head}\label{sec:output}

The final output is produced from the single stream at the end of 
stage $K-1$:
\begin{equation}
  p(x_{t+1} \mid \mathbf{x}_{\leq t}) = 
  \softmax\bigl(W^{\text{out}} \cdot 
  \text{LN}(\mathbf{h}_{t,1}^{(K-1, L_{K-1})}) + 
  \mathbf{b}^{\text{out}}\bigr),
\end{equation}
where $W^{\text{out}} \in \R^{V \times d_{K-1}}$ may optionally be 
tied to the input embedding matrix (transposed).

\subsection{Full Forward Pass}\label{sec:forward}

For completeness, the full forward pass for a sequence $\mathbf{x}$ is 
given in Algorithm~\ref{alg:forward}.

\begin{algorithm}[h]
\caption{Felix-LM Forward Pass}\label{alg:forward}
\begin{algorithmic}[1]
\Require Sequence $\mathbf{x} = (x_1, \ldots, x_T)$; stages 
  $k = 0, \ldots, K{-}1$
\State $\mathbf{e}_t \gets \text{Embed}(x_t)$ for all $t$
\For{$s = 1, \ldots, S_0$} \Comment{Initialize streams}
  \State $\mathbf{h}_{t,s}^{(0,0)} \gets W_s^{\text{init}} 
    \mathbf{e}_t + \mathbf{b}_s^{\text{init}}$ for all $t$
\EndFor
\For{$k = 0, \ldots, K{-}1$} \Comment{Process stages}
  \For{$s = 1, \ldots, S_k$} \Comment{Independent streams}
    \For{$\ell = 0, \ldots, L_k - 1$} \Comment{Layers within stream}
      \State $\mathbf{h}_{:,s}^{(k,\ell+1)} \gets 
        \text{TransformerLayer}_{k,\ell,s}\bigl(
        \mathbf{h}_{:,s}^{(k,\ell)}\bigr)$ 
        \Comment{Eq.~\ref{eq:attn-residual}--\ref{eq:ffn-residual}}
    \EndFor
  \EndFor
  \If{$k < K{-}1$} \Comment{Merge into next stage}
    \State Pair streams: $(s_1, s_2), (s_3, s_4), \ldots$
    \For{each pair $(s, s')$}
      \State $\mathbf{h}_{:,\text{new}}^{(k+1, 0)} \gets 
        \text{Merge}_k\bigl(\mathbf{h}_{:,s}^{(k,L_k)},\; 
        \mathbf{h}_{:,s'}^{(k,L_k)}\bigr)$ 
        \Comment{Eq.~\ref{eq:merge-full}}
    \EndFor
    \State Compute exit loss $\Lcal_k$ 
      \Comment{Eq.~\ref{eq:loss}}
  \EndIf
\EndFor
\State \Return $\softmax\bigl(W^{\text{out}} \cdot 
  \text{LN}(\mathbf{h}_{:,1}^{(K-1, L_{K-1})})\bigr)$
\end{algorithmic}
\end{algorithm}

% ═════════════════════════════════════════════════════════════════════════
\section{Theoretical Analysis}\label{sec:theory}

\subsection{Information-Theoretic Capacity}\label{sec:info-theory}

We analyze the representational capacity of the multi-stream architecture 
relative to a uniform transformer of equal parameter count.

\begin{proposition}[Parameter distribution]\label{prop:params}
Consider a Felix-LM model with $K$ stages, stream counts 
$(S_0, \ldots, S_{K-1})$, dimensions $(d_0, \ldots, d_{K-1})$, and layer 
counts $(L_0, \ldots, L_{K-1})$. Ignoring embedding and output layers, 
the total parameter count is approximately:
\begin{equation}\label{eq:param-count}
  N_{\text{Felix}} \approx \sum_{k=0}^{K-1} S_k \cdot L_k \cdot 
  \bigl(4d_k^2 + 8d_k^2\bigr) + \sum_{k=0}^{K-2} 
  \bigl(4d_k \cdot d_{k+1} + 8d_k^2\bigr),
\end{equation}
where $4d_k^2$ accounts for attention projections ($W_Q, W_K, W_V, W_O$, 
each $d_k \times d_k$, simplified for single-head), $8d_k^2$ for the FFN 
(two layers of $d_k \times 4d_k$), and the merge terms account for the 
gating and projection matrices.

A uniform transformer with $L$ layers of dimension $d$ has:
\begin{equation}
  N_{\text{uniform}} \approx L \cdot 12d^2.
\end{equation}
\end{proposition}

The design constraint is $N_{\text{Felix}} \approx N_{\text{uniform}}$ for 
fair comparison. Given this constraint, the Felix-LM model distributes 
parameters differently: more parameters are allocated to early-stage 
parallel streams (broad exploration) and fewer to late-stage computation 
(focused synthesis, but only one stream).

\begin{proposition}[Representation diversity bound]\label{prop:diversity}
At the output of Stage~0, the $S_0$ independent streams produce $S_0$ 
distinct representations for each token. Under the \textsc{Independent} 
weight configuration, the total representational capacity (measured as the 
dimension of the space of possible token representations) is:
\begin{equation}
  \dim\bigl(\text{range of Stage 0}\bigr) = S_0 \cdot d_0,
\end{equation}
which exceeds the capacity $d_0$ of any individual stream. The merge 
operation projects this $S_0 \cdot d_0$-dimensional representation down to 
$d_1$-dimensional, with the gate $\mathbf{g}_k$ performing learned 
dimensionality selection.

For a uniform transformer, the representation at any intermediate layer 
has capacity $d_{\text{model}}$. If $S_0 \cdot d_0 > d_{\text{model}}$ 
(which holds whenever the streams collectively span more dimensions than 
the uniform model's width), Stage~0 of Felix-LM can represent a strictly 
richer space of initial features before compression.
\end{proposition}

\subsection{Gradient Flow Analysis}\label{sec:gradients}

A critical concern with dimension-changing architectures is gradient flow 
degradation. We analyze the gradient of the loss with respect to early-stage 
parameters.

\begin{proposition}[Gradient path through merges]
\label{prop:grad-merge}
Let $\Lcal$ be the training loss and consider the gradient with respect to 
a parameter $\theta$ in stream $s$ of Stage~0. The gradient flows 
backward through:
\begin{equation}
  \frac{\partial \Lcal}{\partial \theta} = 
  \frac{\partial \Lcal}{\partial \mathbf{h}^{(K-1)}} \cdot
  \prod_{k=K-2}^{0} \underbrace{\frac{\partial 
  \mathbf{h}^{(k+1)}}{\partial \mathbf{h}^{(k)}}}_{\text{merge Jacobian}} 
  \cdot \frac{\partial \mathbf{h}^{(0)}}{\partial \theta}.
\end{equation}

The merge Jacobian has the form:
\begin{equation}\label{eq:merge-jacobian}
  J_k^{\text{merge}} = P_k \cdot \diag\bigl(\mathbf{g}_k + 
  \mathbf{g}_k' \odot [\mathbf{h}_s; \mathbf{h}_{s'}]\bigr),
\end{equation}
where $\mathbf{g}_k' = \mathbf{g}_k \odot (1 - \mathbf{g}_k)$ is the 
sigmoid derivative. If the gate saturates ($g_{k,i} \to 0$ or $g_{k,i} 
\to 1$), the Jacobian for that dimension reduces to $P_k \cdot 
\diag(\mathbf{g}_k)$ (the derivative term vanishes), and gradient flows 
through the non-gated pathway.
\end{proposition}

\begin{corollary}[Gradient preservation under orthogonal merges]
If $P_k$ has orthonormal rows ($P_k P_k^\top = I$), then the merge 
preserves gradient norms in the following sense:
\begin{equation}
  \bigl\|J_k^{\text{merge}} \mathbf{v}\bigr\| \geq 
  \min_i(g_{k,i}) \cdot \|\mathbf{v}\|
\end{equation}
for any gradient vector $\mathbf{v}$, provided all gates are open 
($g_{k,i} = 1$). Gradient vanishing through the merge is thus controlled 
by the gate values, not by the projection geometry.
\end{corollary}

This motivates initializing the gate biases $\mathbf{b}_k^g$ to positive 
values (gates start open) and optionally constraining $P_k$ to be 
orthogonal, especially in deeper models.

\paragraph{Deep supervision as gradient highway.}
The deep supervision loss (\cref{eq:loss}) provides an additional gradient 
pathway: each exit head injects gradients directly into its corresponding 
stage without passing through downstream merges. This ensures that even if 
merge gradients attenuate, every stage receives direct supervisory signal.

\subsection{Representational Geometry}\label{sec:rep-geom}

We now establish that the Felix-LM architecture induces the helical funnel 
trajectory of \cref{def:helical-traj}.

\begin{theorem}[Helical funnel trajectory]\label{thm:helical}
Under the Felix-LM architecture with depth-extended RoPE 
(\cref{def:depth-rope}) and gated merge operations 
(\cref{def:merge}), the representation trajectory 
$\gamma(\mathbf{x})$ satisfies \cref{def:helical-traj} 
with:
\begin{enumerate}[label=(\alph*)]
  \item Angular velocity determined by the depth-RoPE frequencies 
    $\theta_i^{\text{depth}}$ within stages and by the polar decomposition 
    of $P_k$ at merge boundaries;
  \item Radial compression rate $\rho^{(\ell)} = 1$ within stages 
    (constant dimensionality) and $\rho^{(k \to k+1)} = d_{k+1} / 
    (S_k \cdot d_k)$ at merges;
  \item Effective dimensionality decreasing monotonically as 
    $d_{\text{eff}}^{(\text{after stage } k)} = d_{k+1}$, from an initial 
    capacity of $S_0 \cdot d_0$ to a final capacity of $d_{K-1}$.
\end{enumerate}
\end{theorem}
\begin{proof}[Proof sketch]
Within a stage, RMSNorm preserves the subspace, and depth-extended RoPE 
applies an explicit rotation via the depth-frequency term in 
\cref{eq:composite-angle}. This satisfies the 
$Q^{(\ell)}$ term of \cref{eq:helical-decomp} with the gate 
being the identity ($\mathbf{g} = \mathbf{1}$, no compression within 
stages).

At a merge boundary, the concatenation of $S_k$ stream representations 
embeds them in $\R^{S_k \cdot d_k}$. The gated projection 
$P_k(\mathbf{g}_k \odot [\cdot])$ decomposes via polar decomposition 
(\cref{eq:polar}) into rotation ($V_k^\top$, $U_k$) and 
scaling ($\Sigma_k$), with the gate $\mathbf{g}_k$ providing additional 
soft dimensionality reduction. The post-merge dimension $d_{k+1} < S_k 
\cdot d_k$ gives $\rho < 1$, completing the compression step.

The trajectory therefore alternates between rotation-only phases 
(within-stage, via depth-RoPE) and rotation-with-compression phases 
(at merges), tracing a helical funnel as defined.
\end{proof}

% ═════════════════════════════════════════════════════════════════════════
\section{Connection to the Felix Framework}\label{sec:felix-mapping}

The architecture is parameterized to allow direct configuration from the 
Felix framework's YAML specification. The mapping is:

\begin{center}
\renewcommand{\arraystretch}{1.3}
\begin{tabular}{@{}lll@{}}
\toprule
\textbf{Felix Parameter} & \textbf{Felix-LM Analog} & \textbf{Mapping} \\
\midrule
\texttt{top\_radius} ($R_{\max}$) & Initial stream count $S_0$ & 
  $S_0 = 2^{\lceil \log_2 R_{\max} \rceil}$ \\
\texttt{bottom\_radius} ($R_{\min}$) & Final stream count $S_{K-1}$ & 
  $S_{K-1} = 1$ (always) \\
\texttt{turns} ($n$) & Stage transitions $K - 1$ & 
  $K = n + 1$ \\
\texttt{height} ($H$) & Total layers $L_{\text{total}}$ & 
  $L_{\text{total}} \propto H$ \\
\texttt{confidence\_threshold} & Exit threshold $\tau_{\text{conf}}$ & 
  Direct mapping \\
\texttt{temperature\_range} & Attention type per stage & 
  See \S\ref{sec:intra-stream} \\
\texttt{compression\_ratio} & Merge compression & 
  $d_{k+1} / (2 d_k)$ \\
\bottomrule
\end{tabular}
\end{center}

The exponential tapering function (\cref{eq:felix-radius}) determines the 
rate of stream merging: if $R(t)$ decreases exponentially, stream counts 
decrease by factors of 2 at each stage (binary merging), yielding $S_k = 
S_0 / 2^k$. Alternative merging schedules (e.g., merging 3-to-1) 
correspond to steeper radius decay.

% ═════════════════════════════════════════════════════════════════════════
\section{Proposed Experiments}\label{sec:experiments}

\subsection{Ablation Study Design}

All models are trained on WikiText-103 \cite{merity2017pointer} with 
identical optimization hyperparameters. The evaluation metric is 
perplexity on the held-out test set.

\begin{center}
\renewcommand{\arraystretch}{1.3}
\begin{tabular}{@{}clp{8cm}@{}}
\toprule
\textbf{Model} & \textbf{Name} & \textbf{Configuration} \\
\midrule
M0 & \textsc{Uniform} & Standard transformer, $L$ layers, $d$ width. 
  Parameter-matched baseline. \\
M1 & \textsc{MSPM-Homo} & Multi-stream progressive merging; all stages 
  use standard causal attention. Tests the stream-merge structure alone. \\
M2 & \textsc{MSPM-Hetero} & As M1, but Stage~0 uses linear attention, 
  intermediate stages use sliding-window, final stage uses full causal 
  attention. Tests heterogeneous computation. \\
M3 & \textsc{MSPM-Exit} & As M2, plus early-exit heads with 
  cross-stream agreement gating. Tests adaptive depth. \\
M4 & \textsc{Narrow-$d$} & Single stream throughout, but $d_{\text{model}}$ 
  decreases across stages (the ``narrow funnel'' hypothesis). Controls 
  for stream merging vs.\ dimension narrowing. \\
\bottomrule
\end{tabular}
\end{center}

\paragraph{Controlled variables.}
Total parameter count $N \approx 10\text{M}$ ($\pm 5\%$), context length 
$T = 512$, vocabulary from WikiText-103 BPE tokenizer, AdamW optimizer 
with cosine learning rate schedule, batch size, warmup steps, and total 
training tokens are held constant across all models.

\subsection{Hypotheses}

\begin{enumerate}[label=\textbf{H\arabic*.}, leftmargin=2.5em]
  \item M1 achieves lower perplexity than M0, indicating that the 
    multi-stream progressive merging structure provides a useful inductive 
    bias for language modeling.
  \item M2 achieves lower perplexity than M1, indicating that 
    heterogeneous attention types aligned to stage roles improve 
    performance beyond the structural bias alone.
  \item M3 achieves comparable or better perplexity to M2 while reducing 
    average inference FLOPs (via early exit on easy tokens).
  \item M4 achieves worse perplexity than M1, indicating that the funnel 
    should operate on stream count rather than model dimension.
\end{enumerate}

\subsection{Diagnostic Metrics}

Beyond perplexity, we propose monitoring:
\begin{itemize}[nosep]
  \item \textbf{Cross-stream agreement} (\cref{eq:agreement}) over 
    training---do streams converge to specialization or redundancy?
  \item \textbf{Gate activation statistics}: distribution of 
    $\mathbf{g}_k$ values at each merge, tracking which dimensions are 
    preserved versus suppressed.
  \item \textbf{Exit rate by stage}: at inference time, what fraction of 
    tokens exit at each stage?
  \item \textbf{Per-stage loss}: $\Lcal_k$ for each stage, measuring 
    whether early stages learn useful intermediate predictions.
  \item \textbf{Effective dimensionality} (\cref{def:eff-dim}) at each 
    merge output, measuring realized compression.
\end{itemize}

% ═════════════════════════════════════════════════════════════════════════
\section{Extensions and Open Questions}\label{sec:extensions}

\paragraph{Sequence-dimension funneling.}
The current design funnels along the stream/feature axis. An orthogonal 
extension is to pool tokens at merge boundaries (reducing sequence length), 
providing a second axis of compression. The interaction between 
feature-funneling and sequence-funneling is unexplored and theoretically 
interesting: it corresponds to the Felix helix narrowing in both radius 
(feature) and height per turn (sequence).

\paragraph{Skip connections across stages.}
The depth-extended RoPE provides a natural framework for cross-stage 
residual connections, since representations at different depths are in 
well-defined rotational relationships. A skip connection from Stage~0 to 
Stage~2 would require projecting across different dimensions, but the 
rotary structure provides a principled way to align the coordinate frames.

\paragraph{Scaling laws.}
If the architecture proves competitive at $10$M parameters, the next 
question is whether the multi-stream structure obeys favorable scaling 
laws. The additional parameters in merge operations and multiple stream 
initializations represent overhead; whether this overhead amortizes 
favorably with scale is an empirical question.

\paragraph{Hybrid SSM-Attention.}
The heterogeneous attention strategy (\S\ref{sec:intra-stream}) points 
toward a principled framework for combining state-space models (Mamba, 
RWKV) with attention: use SSMs for cheap broad context gathering in early 
stages and attention for precise synthesis in late stages. This is 
consonant with the Felix principle that early agents are broad/cheap and 
late agents are focused/expensive.

\paragraph{Non-binary merging.}
The current design merges streams in pairs. For $S_0$ that is not a power 
of 2, or for more aggressive funneling, $k$-way merges (concatenating $k > 
2$ streams before projection) may be preferable. The gating and projection 
mathematics extend naturally.

% ═════════════════════════════════════════════════════════════════════════
\section{Related Work}\label{sec:related}

\paragraph{Funnel architectures.}
FunnelTransformer \cite{dai2020funnel} narrows the sequence dimension 
via pooling in encoder models. Felix-LM differs in two ways: it narrows 
the stream/feature axis (not sequence), and it targets autoregressive 
(causal) language modeling rather than bidirectional encoding.

\paragraph{Mixture of Experts / Mixture of Depths.}
Mixture-of-Experts models \cite{shazeer2017outrageously, fedus2022switch} 
route tokens to different sub-networks at each layer. Felix-LM's 
multi-stream structure is superficially similar but structurally distinct: 
in MoE, routing is per-token and per-layer; in Felix-LM, all tokens follow 
the same multi-stream structure, and the streams interact only at merge 
points. Mixture of Depths \cite{raposo2024mixture} allows tokens to skip 
layers, which is related to our early-exit mechanism.

\paragraph{Rotary embeddings.}
RoPE \cite{su2021roformer} and its extensions (YaRN 
\cite{peng2023yarn}, NTK-aware scaling) manipulate the rotation 
frequencies for context extension. Our depth-extended RoPE adds a new 
axis (layer index) to the rotation, which to our knowledge has not been 
explored in the literature.

\paragraph{Adaptive computation.}
Universal Transformer \cite{dehghani2019universal} and PonderNet 
\cite{banino2021pondernet} implement depth-adaptive computation via 
halting mechanisms. Felix-LM's cross-stream agreement provides a novel 
confidence signal unavailable in single-stream architectures.

\paragraph{Geometric deep learning.}
Clifford Algebra Networks \cite{ruhe2023geometric} and related work on 
rotation-equivariant architectures \cite{thomas2018tensor, 
fuchs2020se3transformers} provide rigorous frameworks for incorporating 
geometric structure. Felix-LM's helical trajectory is a softer form of 
geometric inductive bias: it does not enforce strict equivariance but 
encourages helical representation dynamics through architectural structure 
and initialization.

% ═════════════════════════════════════════════════════════════════════════
\section{Discussion}\label{sec:discussion}

The central hypothesis of this work is that the principles underlying the 
Felix multi-agent framework---parallel exploration, progressive 
convergence, depth-dependent computation, and confidence-gated resource 
allocation---provide a useful inductive bias when instantiated in the 
architecture of a causal language model.

We have been deliberately conservative in our claims. The proposed 
architecture introduces complexity: multiple streams increase memory 
requirements (though not FLOPs, since streams can be processed 
sequentially), merge operations add parameters, and heterogeneous attention 
types complicate implementation. Whether the inductive bias is strong 
enough to compensate for this added complexity is an empirical question 
that the proposed ablation study is designed to answer.

The most novel contribution, independent of whether the full architecture 
proves competitive, is the \emph{cross-stream agreement} signal for 
confidence estimation. This provides a measure of representational 
robustness---the degree to which independent processing pathways converge 
on the same answer---that has no analog in single-stream architectures 
and may find applications in uncertainty quantification and hallucination 
detection beyond language modeling.

% ═════════════════════════════════════════════════════════════════════════
\section*{Acknowledgments}

The Felix multi-agent framework was developed by Caleb [Surname]. The 
Multi-Stream Progressive Merging architecture was developed through 
iterative collaboration between C.\ [Surname] and AI design assistants 
(Claude, Anthropic). The authors thank [physicist and mathematician 
reviewers] for their feedback on the mathematical formalism.

% ═════════════════════════════════════════════════════════════════════════
% References (manual for now; switch to BibTeX for final version)
\begin{thebibliography}{99}

\bibitem{vaswani2017attention}
A.~Vaswani et al., ``Attention is All You Need,'' \textit{NeurIPS}, 2017.

\bibitem{radford2019language}
A.~Radford et al., ``Language Models are Unsupervised Multitask Learners,'' 
\textit{OpenAI Technical Report}, 2019.

\bibitem{tenney2019bert}
I.~Tenney, D.~Das, E.~Pavlick, ``BERT Rediscovers the Classical NLP 
Pipeline,'' \textit{ACL}, 2019.

\bibitem{clark2019what}
K.~Clark, U.~Khandelwal, O.~Levy, C.D.~Manning, ``What Does BERT Look At? 
An Analysis of BERT's Attention,'' \textit{BlackboxNLP Workshop}, 2019.

\bibitem{su2021roformer}
J.~Su et al., ``RoFormer: Enhanced Transformer with Rotary Position 
Embedding,'' \textit{arXiv:2104.09864}, 2021.

\bibitem{zhang2019root}
B.~Zhang, R.~Sennrich, ``Root Mean Square Layer Normalization,'' 
\textit{NeurIPS}, 2019.

\bibitem{merity2017pointer}
S.~Merity, C.~Xiong, J.~Bradbury, R.~Socher, ``Pointer Sentinel Mixture 
Models,'' \textit{ICLR}, 2017.

\bibitem{dai2020funnel}
Z.~Dai, et al., ``Funnel-Transformer: Filtering out Sequential Redundancy 
for Efficient Language Processing,'' \textit{NeurIPS}, 2020.

\bibitem{shazeer2017outrageously}
N.~Shazeer et al., ``Outrageously Large Neural Networks: The 
Sparsely-Gated Mixture-of-Experts Layer,'' \textit{ICLR}, 2017.

\bibitem{fedus2022switch}
W.~Fedus, B.~Zoph, N.~Shazeer, ``Switch Transformers: Scaling to Trillion 
Parameter Models with Simple and Efficient Sparsity,'' \textit{JMLR}, 2022.

\bibitem{raposo2024mixture}
D.~Raposo et al., ``Mixture-of-Depths: Dynamically Allocating Compute in 
Transformer-Based Language Models,'' \textit{arXiv:2404.02258}, 2024.

\bibitem{peng2023yarn}
B.~Peng et al., ``YaRN: Efficient Context Window Extension of Large 
Language Models,'' \textit{arXiv:2309.00071}, 2023.

\bibitem{dehghani2019universal}
M.~Dehghani et al., ``Universal Transformers,'' \textit{ICLR}, 2019.

\bibitem{banino2021pondernet}
A.~Banino et al., ``PonderNet: Learning to Ponder,'' 
\textit{arXiv:2107.05407}, 2021.

\bibitem{ruhe2023geometric}
D.~Ruhe et al., ``Geometric Clifford Algebra Networks,'' 
\textit{ICML}, 2023.

\bibitem{thomas2018tensor}
N.~Thomas et al., ``Tensor Field Networks: Rotation- and 
Translation-Equivariant Neural Networks for 3D Point Clouds,'' 
\textit{arXiv:1802.08219}, 2018.

\bibitem{fuchs2020se3transformers}
F.~Fuchs et al., ``SE(3)-Transformers: 3D Roto-Translation Equivariant 
Attention Networks,'' \textit{NeurIPS}, 2020.

\end{thebibliography}

% ═════════════════════════════════════════════════════════════════════════
\appendix

\section{Initialization Strategy}\label{sec:init}

The stream initialization projections $W_s^{\text{init}}$ 
(\cref{eq:stream-init}) should encourage diverse initial representations. 
We propose initializing them as orthogonal subspace projections.

Let $d_{\text{embed}} \geq S_0 \cdot d_0$. Generate a random orthogonal 
matrix $Q \in \R^{d_{\text{embed}} \times d_{\text{embed}}}$ (via QR 
decomposition of a random Gaussian matrix) and partition its rows into 
$S_0$ blocks of size $d_0$:
\begin{equation}
  W_s^{\text{init}} = Q_{[(s-1)d_0 : s \cdot d_0,\; :]}.
\end{equation}

This ensures that the initial stream representations are projections onto 
orthogonal subspaces of the embedding, maximizing initial diversity. Under 
this initialization:
\begin{equation}
  \langle W_s^{\text{init}} \mathbf{e}, W_{s'}^{\text{init}} 
  \mathbf{e}\rangle = 0 \quad \text{for } s \neq s',
\end{equation}
for any input $\mathbf{e}$, guaranteeing that streams start with 
completely independent views of the input.

\section{Parameter Count Calculations}\label{sec:param-calc}

We provide exact parameter count formulas for the ablation models, 
parameterized to target $N \approx 10$M.

For each model, the total parameters are:
\[
  N = N_{\text{embed}} + N_{\text{stages}} + N_{\text{merges}} + 
  N_{\text{exits}} + N_{\text{output}},
\]
where:
\begin{align}
  N_{\text{embed}} &= V \cdot d_{\text{embed}} + S_0 \cdot d_0 \cdot 
    d_{\text{embed}}, \\
  N_{\text{stages}} &= \sum_{k=0}^{K-1} S_k \cdot L_k \cdot 
    \bigl(4 d_k^2 + 2 \cdot d_k \cdot 4d_k\bigr), \\
  N_{\text{merges}} &= \sum_{k=0}^{K-2} \frac{S_k}{2} \cdot 
    \bigl(d_{k+1} \cdot 2d_k + 2d_k \cdot 2d_k\bigr), \\
  N_{\text{exits}} &= \sum_{k=0}^{K-2} \bigl(d_k \cdot V\bigr), \\
  N_{\text{output}} &= d_{K-1} \cdot V.
\end{align}

Concrete configurations for each model will be computed numerically under 
the constraint $N = 10 \times 10^6 \pm 5\%$.

\end{document}
